\documentclass[11pt,a4paper]{article}
% Preambulo comun de los reportes de Tecnologias Celulares.
% Cada main.tex hace: \documentclass ... % Preambulo comun de los reportes de Tecnologias Celulares.
% Cada main.tex hace: \documentclass ... % Preambulo comun de los reportes de Tecnologias Celulares.
% Cada main.tex hace: \documentclass ... \input{../comun/preambulo.tex}
\usepackage[utf8]{inputenc}
\usepackage[T1]{fontenc}
\usepackage[spanish,mexico]{babel}
\usepackage{graphicx}
\usepackage[margin=2.5cm]{geometry}
\usepackage{amsmath}
\usepackage{amssymb}
\usepackage{booktabs}
\usepackage{tabularx}
\usepackage{array}
\usepackage{float}
\usepackage{setspace}
\usepackage{parskip}
\usepackage{fancyhdr}
\usepackage{textcomp}
\usepackage{xcolor}
\usepackage{caption}
\usepackage{fvextra}
\usepackage[hidelinks]{hyperref}

\onehalfspacing
\setlength{\parskip}{6pt}
\sloppy
\setlength{\emergencystretch}{3em}
\setlength{\headheight}{40pt}
\captionsetup{font=small,labelfont=bf}
\setlength{\tabcolsep}{6pt}

\pagestyle{fancy}
\fancyhf{}
\fancyhead[L]{\includegraphics[width=1cm]{/home/yovick/repos/UAZ-IEI/Assets/UIE.png}}
\fancyhead[R]{\small Tecnologías Celulares}
\fancyfoot[C]{\thepage}
\renewcommand{\headrulewidth}{0.4pt}

\DefineVerbatimEnvironment{codigo}{Verbatim}{%
  breaklines=true,breakanywhere=true,fontsize=\small,frame=single,
  framesep=3pt,rulecolor=\color{gray!50}}

% Figura centrada con pie: \figura[ancho]{archivo}{leyenda}
\newcommand{\figura}[3][0.85]{%
  \begin{center}
    \includegraphics[width=#1\linewidth]{#2}\\[2pt]
    {\small #3}
  \end{center}}

% Portada: \portada{numero}{titulo de la practica}
\newcommand{\portada}[2]{%
\begin{titlepage}
  \thispagestyle{empty}
  \begin{center}
    \vspace*{1cm}
    \includegraphics[width=4.5cm]{/home/yovick/repos/UAZ-IEI/Assets/uaz.png}
    \vspace{0.8cm}

    {\large Universidad Autónoma de Zacatecas}\\[4pt]
    {\large Unidad Académica de Ingeniería Eléctrica}\\[4pt]
    {\normalsize Ingeniería Electrónica Industrial con especialidad en Telecomunicaciones}

    \vspace{1.8cm}
    {\LARGE\bfseries Práctica #1}\\[10pt]
    {\Large\bfseries #2}
  \end{center}

  \vspace{2.2cm}
  \begin{center}
    \begin{tabular}{rl}
      \textbf{Materia:} & Tecnologías Celulares\\
      \textbf{Docente:} & Dr. Víktor Iván Rodríguez Abdalá\\
      \textbf{Alumno:}  & Arnold Yovick Ríos Zaldívar\\
      \textbf{Fecha:}   & Noviembre de 2026\\
    \end{tabular}
  \end{center}
  \vfill
\end{titlepage}}

\usepackage[utf8]{inputenc}
\usepackage[T1]{fontenc}
\usepackage[spanish,mexico]{babel}
\usepackage{graphicx}
\usepackage[margin=2.5cm]{geometry}
\usepackage{amsmath}
\usepackage{amssymb}
\usepackage{booktabs}
\usepackage{tabularx}
\usepackage{array}
\usepackage{float}
\usepackage{setspace}
\usepackage{parskip}
\usepackage{fancyhdr}
\usepackage{textcomp}
\usepackage{xcolor}
\usepackage{caption}
\usepackage{fvextra}
\usepackage[hidelinks]{hyperref}

\onehalfspacing
\setlength{\parskip}{6pt}
\sloppy
\setlength{\emergencystretch}{3em}
\setlength{\headheight}{40pt}
\captionsetup{font=small,labelfont=bf}
\setlength{\tabcolsep}{6pt}

\pagestyle{fancy}
\fancyhf{}
\fancyhead[L]{\includegraphics[width=1cm]{/home/yovick/repos/UAZ-IEI/Assets/UIE.png}}
\fancyhead[R]{\small Tecnologías Celulares}
\fancyfoot[C]{\thepage}
\renewcommand{\headrulewidth}{0.4pt}

\DefineVerbatimEnvironment{codigo}{Verbatim}{%
  breaklines=true,breakanywhere=true,fontsize=\small,frame=single,
  framesep=3pt,rulecolor=\color{gray!50}}

% Figura centrada con pie: \figura[ancho]{archivo}{leyenda}
\newcommand{\figura}[3][0.85]{%
  \begin{center}
    \includegraphics[width=#1\linewidth]{#2}\\[2pt]
    {\small #3}
  \end{center}}

% Portada: \portada{numero}{titulo de la practica}
\newcommand{\portada}[2]{%
\begin{titlepage}
  \thispagestyle{empty}
  \begin{center}
    \vspace*{1cm}
    \includegraphics[width=4.5cm]{/home/yovick/repos/UAZ-IEI/Assets/uaz.png}
    \vspace{0.8cm}

    {\large Universidad Autónoma de Zacatecas}\\[4pt]
    {\large Unidad Académica de Ingeniería Eléctrica}\\[4pt]
    {\normalsize Ingeniería Electrónica Industrial con especialidad en Telecomunicaciones}

    \vspace{1.8cm}
    {\LARGE\bfseries Práctica #1}\\[10pt]
    {\Large\bfseries #2}
  \end{center}

  \vspace{2.2cm}
  \begin{center}
    \begin{tabular}{rl}
      \textbf{Materia:} & Tecnologías Celulares\\
      \textbf{Docente:} & Dr. Víktor Iván Rodríguez Abdalá\\
      \textbf{Alumno:}  & Arnold Yovick Ríos Zaldívar\\
      \textbf{Fecha:}   & Noviembre de 2026\\
    \end{tabular}
  \end{center}
  \vfill
\end{titlepage}}

\usepackage[utf8]{inputenc}
\usepackage[T1]{fontenc}
\usepackage[spanish,mexico]{babel}
\usepackage{graphicx}
\usepackage[margin=2.5cm]{geometry}
\usepackage{amsmath}
\usepackage{amssymb}
\usepackage{booktabs}
\usepackage{tabularx}
\usepackage{array}
\usepackage{float}
\usepackage{setspace}
\usepackage{parskip}
\usepackage{fancyhdr}
\usepackage{textcomp}
\usepackage{xcolor}
\usepackage{caption}
\usepackage{fvextra}
\usepackage[hidelinks]{hyperref}

\onehalfspacing
\setlength{\parskip}{6pt}
\sloppy
\setlength{\emergencystretch}{3em}
\setlength{\headheight}{40pt}
\captionsetup{font=small,labelfont=bf}
\setlength{\tabcolsep}{6pt}

\pagestyle{fancy}
\fancyhf{}
\fancyhead[L]{\includegraphics[width=1cm]{/home/yovick/repos/UAZ-IEI/Assets/UIE.png}}
\fancyhead[R]{\small Tecnologías Celulares}
\fancyfoot[C]{\thepage}
\renewcommand{\headrulewidth}{0.4pt}

\DefineVerbatimEnvironment{codigo}{Verbatim}{%
  breaklines=true,breakanywhere=true,fontsize=\small,frame=single,
  framesep=3pt,rulecolor=\color{gray!50}}

% Figura centrada con pie: \figura[ancho]{archivo}{leyenda}
\newcommand{\figura}[3][0.85]{%
  \begin{center}
    \includegraphics[width=#1\linewidth]{#2}\\[2pt]
    {\small #3}
  \end{center}}

% Portada: \portada{numero}{titulo de la practica}
\newcommand{\portada}[2]{%
\begin{titlepage}
  \thispagestyle{empty}
  \begin{center}
    \vspace*{1cm}
    \includegraphics[width=4.5cm]{/home/yovick/repos/UAZ-IEI/Assets/uaz.png}
    \vspace{0.8cm}

    {\large Universidad Autónoma de Zacatecas}\\[4pt]
    {\large Unidad Académica de Ingeniería Eléctrica}\\[4pt]
    {\normalsize Ingeniería Electrónica Industrial con especialidad en Telecomunicaciones}

    \vspace{1.8cm}
    {\LARGE\bfseries Práctica #1}\\[10pt]
    {\Large\bfseries #2}
  \end{center}

  \vspace{2.2cm}
  \begin{center}
    \begin{tabular}{rl}
      \textbf{Materia:} & Tecnologías Celulares\\
      \textbf{Docente:} & Dr. Víktor Iván Rodríguez Abdalá\\
      \textbf{Alumno:}  & Arnold Yovick Ríos Zaldívar\\
      \textbf{Fecha:}   & Noviembre de 2026\\
    \end{tabular}
  \end{center}
  \vfill
\end{titlepage}}


\begin{document}

\portada{4}{Implementación de un sistema FHSS}

\textbf{Script:} \texttt{practica4.py} (Python con NumPy/SciPy)

\section{Objetivo}

Implementar un sistema de espectro disperso por salto de frecuencia (FHSS), generar la
secuencia pseudoaleatoria de saltos y evaluar su desempeño frente a una interferencia de
banda parcial, comparándolo con una transmisión convencional de frecuencia fija.

\section{Fundamento teórico}

En FHSS la portadora cambia de frecuencia en cada símbolo siguiendo una secuencia
pseudoaleatoria conocida por el transmisor y el receptor. Si hay $M$ canales disponibles,
\begin{equation}
  f_h(k) = f_{m(k)},
\end{equation}
donde $m(k)\in\{1,\dots,M\}$ es el índice generado por la secuencia PN. El receptor conoce
la secuencia y sintoniza el mismo canal en cada intervalo.

La ventaja aparece cuando la interferencia ocupa solo una parte del espectro: si el canal
$j$ está interferido, el sistema convencional que transmite siempre en $j$ sufre la
interferencia en el 100\,\% de los símbolos, mientras que el FHSS solo la sufre en la
fracción $1/M$. Para una interferencia gaussiana de desviación $A_I$ sobre BPSK, la
probabilidad de error de un símbolo afectado es $Q(1/A_I)$ y la BER promedio resulta
\begin{equation}
  P_b^{\mathrm{FHSS}} = \frac{1}{M}\,Q\!\left(\frac{1}{A_I}\right),
  \qquad
  P_b^{\mathrm{conv}} = Q\!\left(\frac{1}{A_I}\right).
\end{equation}
El factor $1/M$ es la ganancia de diversidad en frecuencia: el sistema no evita la
interferencia, pero la reparte en el tiempo.

\section{Desarrollo experimental}

\subsection{Parámetros}

\begin{tabularx}{\linewidth}{@{}*{2}{>{\raggedright\arraybackslash}X}@{}}
  \toprule
  Parámetro & Valor \\
  \midrule
  Bits de la demostración & 200 \\
  Canales disponibles & $M=8$ \\
  Canal interferido & 4 \\
  Modulación & BPSK de amplitud unitaria \\
  Interferencia & ruido gaussiano real de desviación $A_I$ \\
  Niveles evaluados & $A_I=2,4,6,8,10$ \\
  \bottomrule
\end{tabularx}

\subsection{Modelo}

El manual multiplica los bits por el número de canal (\texttt{fhSignal = hops.*b}) y luego
añade ruido solo en los símbolos del canal 4. Ese modelo no mide bien la interferencia:
como la amplitud del símbolo crece con el número de canal, la detección por signo casi
nunca falla. En este trabajo se usa el modelo equivalente correcto: cada símbolo BPSK tiene
amplitud unitaria y se transmite en el canal seleccionado; el receptor conoce el salto y
detecta por signo. La interferencia se suma únicamente cuando el salto cae en el canal 4.
El sistema convencional transmite siempre en el canal 4, por lo que recibe la interferencia
en todos los símbolos.

\subsection{Procedimiento}

\begin{enumerate}
  \item Se generan 200 bits y 200 saltos pseudoaleatorios uniformes entre 1 y 8.
  \item Se construye la señal transmitida y se contamina con interferencia en los símbolos
        cuyo salto es el canal 4.
  \item Se repite con 200\,000 bits por nivel de $A_I$ para obtener curvas estables.
  \item Se calcula la BER del FHSS y la de un sistema convencional con portadora fija en
        el canal interferido.
\end{enumerate}

\section{Código}

\begin{codigo}
def simular(ai, nbits=200_000, rng=None):
    bits = rng.integers(0, 2, nbits)
    b = 2 * bits - 1
    hops = rng.integers(1, CANALES + 1, nbits)
    afectados = hops == CANAL_JAM

    rx = b.astype(float).copy()
    rx[afectados] += ai * rng.standard_normal(afectados.sum())
    ber_fhss = np.mean((rx > 0).astype(int) != bits)

    rx_conv = b + ai * rng.standard_normal(nbits)
    ber_conv = np.mean((rx_conv > 0).astype(int) != bits)
    return ber_fhss, ber_conv
\end{codigo}

\section{Resultados}

\subsection{Secuencia de salto}

\figura{figs/saltos.png}{Secuencia de salto para 200 símbolos; en magenta, los que caen en el canal interferido.}

Los saltos se distribuyen de forma aproximadamente uniforme entre los ocho canales. En la
demostración, 21 de los 200 símbolos (10.5\,\%) cayeron en el canal 4, valor cercano al
12.5\,\% esperado.

\subsection{Representación tiempo--frecuencia}

\figura{figs/tiempo_frecuencia.png}{Cada símbolo ocupa un canal distinto; la banda sombreada es el canal interferido.}

\subsection{Señal afectada por la interferencia}

\figura{figs/senal_interferida.png}{Señal transmitida y recibida con $A_I=8$. Solo los símbolos del canal 4 se separan del valor nominal.}

\subsection{BER frente a la interferencia}

\begin{tabular}{@{}cccc@{}}
  \toprule
  $A_I$ & BER FHSS & BER convencional & Teórica FHSS \\
  \midrule
  2  & 0.03851 & 0.30817 & 0.03857 \\
  4  & 0.05018 & 0.40043 & 0.05016 \\
  6  & 0.05422 & 0.43217 & 0.05423 \\
  8  & 0.05715 & 0.45138 & 0.05628 \\
  10 & 0.05795 & 0.46060 & 0.05752 \\
  \bottomrule
\end{tabular}

\figura{figs/ber.png}{Comparación FHSS contra transmisión convencional.}

La BER del FHSS es aproximadamente ocho veces menor que la del sistema convencional en
todos los niveles, que es exactamente la fracción de símbolos afectados. La simulación
coincide con el modelo teórico $Q(1/A_I)/8$ dentro del error de Monte Carlo.

\section{Análisis de resultados}

\begin{itemize}
  \item \textbf{¿Cómo afecta la interferencia a cada sistema?} Al sistema convencional lo
        afecta en todos los símbolos porque su portadora está fija en el canal interferido.
        Al FHSS solo lo afecta cuando el salto cae en ese canal, es decir, una fracción
        $1/M$ del tiempo.
  \item \textbf{¿Por qué FHSS es más robusto?} Porque la diversidad en frecuencia convierte
        una interferencia continua en una interferencia intermitente: el receptor pierde
        pocos símbolos y puede recuperar el resto, o corregirlos con codificación de canal.
  \item \textbf{¿Qué ocurre si se incrementa el número de frecuencias disponibles?} La
        fracción de símbolos afectados baja como $1/M$, así que la BER promedio mejora.
        El costo es que el receptor debe cubrir un ancho de banda mayor y sincronizar
        saltos más rápidos.
  \item \textbf{¿Cómo influye la secuencia de salto en el desempeño?} Una secuencia con
        buena uniformidad y bajo solapamiento evita que el transmisor permanezca mucho
        tiempo en un canal interferido. Si la secuencia es predecible, un interferidor
        inteligente puede seguirla y anular la ventaja.
\end{itemize}

\section{Cuestionario}

\begin{enumerate}
  \item \textbf{¿Qué es FHSS?} Espectro disperso por salto de frecuencia: la portadora
        cambia de frecuencia siguiendo una secuencia pseudoaleatoria, de modo que la
        transmisión ocupa en promedio todo el ancho de banda disponible.
  \item \textbf{¿Cuál es la principal diferencia entre FHSS y DSSS?} DSSS transmite
        simultáneamente en toda la banda y dispersa con un código; FHSS transmite en una
        banda angosta por vez y cambia de frecuencia con el tiempo.
  \item \textbf{¿Qué papel desempeña la secuencia pseudoaleatoria?} Define el orden de los
        saltos. Debe ser conocida por el receptor para sintonizar el canal correcto y debe
        tener buena uniformidad para no concentrarse en canales interferidos.
  \item \textbf{¿Por qué FHSS proporciona diversidad en frecuencia?} Porque la información
        se reparte en muchos canales distintos; si uno está interferido o atenuado, solo
        se pierde la fracción de símbolos que cayó en él.
  \item \textbf{¿Qué aplicaciones comerciales utilizan FHSS?} Bluetooth clásico (79 canales
        y 1600 saltos por segundo), algunos sistemas inalámbricos de telefonía DECT,
        micrófonos y radios militares, y redes de sensores que buscan robustez frente a
        interferencia.
\end{enumerate}

\section{Conclusiones}

Se implementó un sistema FHSS de ocho canales y se comparó con una transmisión de
portadora fija frente a una interferencia localizada en el canal 4. La BER del FHSS varió
entre 0.0385 y 0.0580 para $A_I$ de 2 a 10, mientras que el sistema convencional varió
entre 0.3082 y 0.4606: una reducción de aproximadamente ocho veces, igual a la fracción
$1/M$ de símbolos afectados. El resultado confirma que la diversidad en frecuencia es la
ventaja central del FHSS frente a interferencias de banda parcial.

\end{document}
